% ==============================================================
\section{Experiments}\label{sec:experiments}
% ==============================================================

\subsection{Experimental Design}

\subsubsection{Synthetic Data Generation}

We generate synthetic data with a base vocabulary of 200 items
(per-item occurrence probability $p_{\text{base}} = 0.03$)
and $N = 100{,}000$ transactions.
External event effects are modeled as \emph{planted signals}.
Three signal types are mixed in a single dataset to
multi-dimensionally evaluate the pipeline's discrimination ability:

\noindent\textbf{Type A} (vocab-in boost):
During the event active period, all pattern items are simultaneously inserted
with probability $\beta$, raising support.
A realistic condition in which baseline support exists at $p_{\text{base}}$ outside the event period.

\noindent\textbf{Type B} (event-unrelated dense patterns):
Support rises in a specific period, but not corresponding to any event
(simulates seasonality or trend).

\noindent\textbf{Type C} (decoy events):
Events unrelated to any pattern variation.

\subsubsection{Evaluation Metrics}\label{sec:eval_metrics}

Evaluation is performed at the \textbf{(P, I, E) triplet level}.

\paragraph{Ground Truth Construction}
For each planted signal $(P, E)$, the longest dense interval overlapping
the event window is registered as the representative $I^*$,
yielding ground truth triplet $(P, I^*, E)$ (one per $(P, E)$ pair).

\paragraph{Matching Criterion}
Exact match on item set, interval endpoints, and event ID is required.

\paragraph{Precision Metrics}
Standard Precision, Recall, and F1 score based on
$\text{TP} = |\hat{\mathcal{A}} \cap \mathcal{A}^*|$,
$\text{FP} = |\hat{\mathcal{A}} \setminus \mathcal{A}^*|$,
$\text{FN} = |\mathcal{A}^* \setminus \hat{\mathcal{A}}|$.
Evaluation is performed on the final output after deduplication (Step~5).
Type~B (event-unrelated) patterns output as significant attributions
count as FP and reduce Precision.

\subsubsection{Default Settings}

Default parameters are shown in Appendix~\ref{app:params} (Table~\ref{tab:default_params}).
Square root normalization (Eq.~(\ref{eq:mag_sqrt})) is used for magnitude normalization.

\subsubsection{Evaluation Policy}

EX1 evaluates attribution accuracy under diverse data conditions.
EX2 demonstrates the contribution of each step through ablation analysis.
EX3 performs a step-difference comparison replacing Steps~1, 3, and 4 with related methods.
EX4 confirms real-world data applicability through exploratory analysis;
quantitative evaluation is not performed because ground truth attribution labels
are unavailable for real-world data.

% ------------------------------------------------------------------
\subsection{EX1: Core Attribution Accuracy}
\label{sec:exp_ex1}

\textbf{RQ}: Can the pipeline correctly attribute event-driven variation
and reject event-unrelated dense patterns?
Can it naturally handle multiple pattern lengths $L = 2$--$4$?
We evaluate robustness across diverse data characteristics
(signal strength, structural conditions, pattern length).

\paragraph{Signal Strength and Structural Conditions}
Type A (5 vocab-in boost patterns: $L=2\times3$, $L=3\times1$, $L=4\times1$),
Type B (2 event-unrelated patterns), and Type C (2 decoys) are mixed,
evaluated with \textbf{20 seeds} across $\beta \in \{0.1, 0.2, 0.3, 0.4, 0.5\}$
and 4 structural conditions (Fig.~\ref{fig:ex1_combined}).

\begin{figure}[t]
\centering
\includegraphics[width=\linewidth]{ex1_combined.pdf}
\caption{EX1 results (20-seed mean with 95\% CI).
(a) Signal strength: $\beta \geq 0.2$ achieves F1 $\geq$ 0.78 and Recall $\geq$ 0.92.
$\beta = 0.1$ reaches the detection limit for $L=2$ vocab-in patterns (F1 = 0.34),
while $L=3,4$ vocab-out patterns are stably detected.
(b) Structural conditions: F1 $\geq$ 0.67 across all conditions
(SHORT: 0.85, OVERLAP: 0.74, CONFOUND: 0.70, DENSE: 0.67).}
\label{fig:ex1_combined}
\end{figure}

At $\beta = 0.1$, $L=2$ vocab-in patterns reach the detection limit
while $L=3,4$ vocab-out patterns are stably detected:
F1 = 0.34~[0.30, 0.37] (Recall = 0.22).
At $\beta \geq 0.2$, Recall $\geq 0.92$ and F1 $\geq 0.78$ are achieved.
F1 $\geq$ 0.67 is achieved across all structural conditions
(see \S\ref{app:overlap_analysis}).

\paragraph{Null Condition Validation}
Running \textbf{100 seeds} on null data with no planted signals (Appendix~\ref{app:null_fdr}),
the mean per-seed FDR was \textbf{0.070} ($\leq \alpha = 0.10$), empirically confirming FDR control.

% ------------------------------------------------------------------
\subsection{EX2: Ablation Analysis}
\label{sec:exp_ex2}

\textbf{RQ}: How does each score component and pipeline step contribute to accuracy?
We use Zipf-distributed data ($\sqrt{}$ normalization) as in EX1.

\paragraph{Score Components}

We design two scenarios demonstrating the necessity of each component
and apply all 3 ablation modes to each scenario
(2 scenarios $\times$ 3 modes $\times$ 5 seeds = 30 conditions):

\noindent\textbf{Scenario A} (prox required):
Pattern $(5, 15)$ has two dense intervals---one near the event (causal)
and one distant (accidental). Without prox, both receive equal scores.

\noindent\textbf{Scenario B} (mag required):
Pattern $(5, 15)$ is affected by 2 events with different boost magnitudes
($\beta=0.5$ vs.\ $\beta=0.1$). Without mag, both are treated equally.

Results are shown in Table~\ref{tab:ex2a}.

\begin{table}[t]
\caption{EX2a Results: Per-Scenario Ablation (5-seed mean F1, Zipf-1.0, $\sqrt{}$ normalization)}
\label{tab:ex2a}
\centering
\begin{tabular}{lcc}
\toprule
Configuration & Scenario A & Scenario B \\
\midrule
Full: $\text{prox} \times \text{mag}$ & \textbf{1.00} & \textbf{0.46} \\
mag only (prox removed)               & 0.00 & 0.00 \\
prox only (mag removed)               & 1.00 & 0.19 \\
\bottomrule
\end{tabular}
\end{table}

Without prox (mag only), the method has no temporal alignment information,
collapsing to \textbf{F1 = 0.00 in both scenarios}.
Prox only achieves F1 = 1.00 on Scenario~A but drops to F1 = 0.19
on Scenario~B (distinguishing strong vs.\ weak boost).
The full configuration achieves 1.00 on Scenario~A and 0.46 on Scenario~B,
demonstrating the complementarity of both components.

\paragraph{Pipeline Step Contributions}
We compare a naive baseline (Steps~1--3 only: change point detection + scoring,
outputting all candidates above threshold) against three variants
obtained by incrementally adding pipeline steps:

\begin{itemize}
\item \textbf{Naive}: Steps~1--3 only (no testing)
\item \textbf{+PermTest}: Per-pattern permutation test added (no global correction)
\item \textbf{+BH}: Global BH correction added
\item \textbf{Full}: Union-Find deduplication added (= complete pipeline)
\end{itemize}

Using EX1 data ($\beta = 0.3$, CONFOUND, DENSE; Zipf-1.0, $\sqrt{}$ normalization),
evaluated with 5 seeds per condition (Fig.~\ref{fig:ex2b_ablation}).

\begin{figure}[t]
\centering
\includegraphics[width=\linewidth]{ex2b_ablation.pdf}
\caption{EX2b: Accuracy change from incremental step additions (5-seed mean, Zipf-1.0, $\sqrt{}$ normalization).
(a) F1: substantial improvement from BH correction and deduplication.
(b) \#Predictions: each step progressively reduces output volume, suppressing false positives.}
\label{fig:ex2b_ablation}
\end{figure}

Naive achieves Recall = 1.00 across all conditions but F1 $\leq$ 0.03
due to massive false positives---score thresholding alone is insufficient.
\textbf{+PermTest} provides limited FP reduction,
but \textbf{+BH} (global BH correction) delivers the largest false positive reduction,
and \textbf{Full} (+deduplication) further eliminates redundant patterns
to achieve the best F1 across all conditions.
Each step is confirmed to reduce false positives in a progressive and complementary manner.

% ------------------------------------------------------------------
\subsection{EX3: Comparison with Related Methods}
\label{sec:exp_ex3}

\textbf{RQ}: Is the proposed component design optimal when the pipeline's
Steps~1--3--4 are replaced with related methods?

\noindent\textbf{Comparison Purpose and Design}
Step~0 (Apriori-window), BH correction, and Union-Find deduplication (Step~5)
are shared by all methods; only the attribution core (Steps~1, 3, 4) is replaced.
Results should be interpreted as ``proposed step design vs.\ alternative step designs.''

From the related methods in Table~\ref{tab:related_comparison},
we select methods that test the association between events and outcomes,
replacing Steps~1 (change point extraction), 3 (attribution scoring),
and 4 (statistical testing) with the following four methods:
\begin{enumerate}
\item \textbf{Wilcoxon}~\cite{mann1947test}: Rank-sum test comparing in-event vs.\ out-of-event support distributions.
\item \textbf{CausalImpact}~\cite{brodersen2015causalimpact}: Local-level model estimates counterfactual and tests cumulative deviation.
\item \textbf{ITS}~\cite{bernal2017interrupted}: Segmented regression tests level change coefficient before and after intervention.
\item \textbf{EventStudy}~\cite{mackinlay1997event}: Cumulative abnormal support tested via $z$-test from pre-event statistics.
\end{enumerate}

Results on the same data as EX1 ($\beta = 0.3$, Zipf-1.0, $\sqrt{}$ normalization)
and 4 structural conditions (5 seeds each) are shown in Fig.~\ref{fig:ex3_method}.

\begin{figure}[t]
\centering
\includegraphics[width=\linewidth]{ex3_method.pdf}
\caption{EX3 results: comparison with related methods (5-seed mean F1).
All methods share Step~0 (Apriori-window), BH correction, and deduplication.
The proposed method achieves the highest or equal F1 across all 5 conditions,
with the most prominent advantage on CONFOUND (0.68 vs.\ runner-up ITS 0.45).}
\label{fig:ex3_method}
\end{figure}

The proposed method achieves the highest or equal F1 across all 5 conditions
(near-tied with ITS on SHORT: 0.72 vs.\ 0.77),
particularly exceeding runner-up ITS on CONFOUND (F1 = 0.68 vs.\ 0.45).
The proximity score leveraging change point timing effectively discriminates
event-unrelated patterns whose active periods overlap with event windows.

ITS is competitive at $\beta=0.3$ (F1 = 0.60 vs.\ proposed: 0.66)
but degrades significantly on CONFOUND (0.45) and OVERLAP (0.40).
CausalImpact and Wilcoxon remain at F1 $\leq$ 0.49 across all conditions.
The proposed method achieves F1 0.66--0.72 (5-condition average 0.68),
with consistent stability and marked advantage under confounding and overlap conditions.

% ------------------------------------------------------------------
\subsection{EX4: Real-World Data Validation}
\label{sec:exp_ex4}

\textbf{RQ}: Is the method applicable to real-world data?

Exploratory analysis on 276{,}484 purchase baskets from
the Dunnhumby ``The Complete Journey'' dataset~\cite{dunnhumby2014complete}.
Square root normalization is used, as item frequencies follow a Zipf-like distribution.

\paragraph{External Event Selection}
This pipeline targets population-level support variation (Definition~\ref{def:attribution}).
The Dunnhumby dataset records 30 promotional campaigns,
but targeted campaigns (TypeB/C: fewer than 2{,}000 coupon items)
target specific households with individual coupons
and are not expected to produce detectable population-level effects.
We therefore use the \textbf{5 large-scale promotional campaigns}
(C8, C13, C18, C26, C30; 3.7K--38.2K coupon items)
with at least 3{,}000 coupon items as external events.
This selection is based on campaign metadata (coupon.csv) as a pre-hoc filter,
not selected post-hoc by consulting attribution results.

Settings: $W = 300$, $\theta = 5$, $B = 5{,}000$, $\alpha = 0.10$ (BH correction),
deduplication enabled.

\begin{table}[t]
\caption{EX4 Results: Dunnhumby Large-Scale Promotional Campaigns}
\label{tab:ex4_dunnhumby}
\centering
\begin{tabular}{lr}
\toprule
Metric & Value \\
\midrule
Transactions & 276{,}484 \\
Patterns ($|P| \geq 2$) & 896 \\
Campaigns & 5 \\
Significant attributions (post-dedup) & 15 \\
Mean attributions per campaign & 3.0 \\
\bottomrule
\end{tabular}
\end{table}

At least one significant attribution was generated for all 5 campaigns
(100\% coverage), with an average of 3.0 attributions per campaign---a
manageable volume for analysts.
Quantitative accuracy evaluation is not performed because ground truth labels are unavailable,
but the product categories of attributed patterns qualitatively align with campaign content,
confirming applicability to real-world data.
